\section{Introduction}

Recent advances in transformer-based large language models have demonstrated remarkable capabilities in financial domains, from sentiment extraction to return forecasting~\cite{lopez2023can,chen2023fingpt}. Yet a fundamental question persists: Do these systems genuinely comprehend market microstructure mechanisms, or are they sophisticated pattern-matching engines that exploit memorized associations from training corpora?

This distinction carries profound implications. Options market makers operate under regulatory mandates to maintain delta-neutral books despite accumulating gamma exposure~\cite{sec15c31}. When aggregate gamma exposure turns substantially negative, dealers face forced pro-cyclical hedging—selling strength and purchasing weakness—thereby amplifying price movements through coordinated rebalancing flows~\cite{ni2005stock,garleanu2009demand}. These market-wide patterns emerge not from centralized coordination but from countless dealers independently responding to local constraints—what Hayek termed emergent order arising from dispersed knowledge~\cite{hayek1945}. Validating whether LLMs detect such structural constraints presents a methodological challenge: presenting a model with ``SPY on January 27, 2021'' alongside substantial negative gamma risks measuring memorization of the GameStop squeeze rather than comprehension of underlying dealer mechanics. This training data contamination problem pervades LLM evaluation in finance, where temporal context enables recall rather than reasoning.

We address this through \textit{obfuscation testing}, a validation approach that eliminates all memorizable information while preserving structural relationships. Converting dates to generic labels (``Day T+0''), anonymizing ticker symbols (``SPY'' → ``INDEX\_1''), and removing economic event references forces the LLM to reason exclusively from market structure—gamma exposure magnitudes, strike distributions, open interest concentrations. Successful pattern detection under these conditions demonstrates understanding of causal constraint mechanisms rather than statistical correlation with training data.

The proliferation of zero-days-to-expiration (0DTE) options establishes a particularly rigorous validation environment. Trading volume in 0DTE contracts expanded from negligible levels pre-2022 to exceeding 50\% of total SPX options volume by 2024~\cite{cboe2023}, generating a persistent negative gamma regime where dealers maintain net short gamma across 95.6\% of trading days in our sample. This structural transition from historical alternating regimes~\cite{ni2005stock,garleanu2009demand} to sustained single-regime dynamics presents a demanding test case: unlike conventional pattern recognition exploiting regime variation, our methodology must identify the underlying constraint mechanism within a uniform environment where hedging dynamics evolve from amplification-dominant (early 2024) to equilibrium-dominant (late 2024) as market participants adapt to persistent dealer short positioning.

Our methodology advances LLM validation in financial domains through four contributions. First, we introduce obfuscation testing for distinguishing genuine structural comprehension from memorization, establishing a rigorous framework for domains where training data contamination threatens validity. Second, we quantify prompt bias effects on detection sensitivity: including regime labels (``NEGATIVE\_GAMMA'') inflates detection from 71.5\% to 100\%, revealing how contextual hints generate misleading performance metrics. Third, we validate detection robustness across multiple pattern framings and temporal periods, demonstrating generalization rather than overfitting to specific narratives. Fourth, we establish detection-profitability divergence: stable detection capability (68-74\% quarterly) persists as economic alpha declines to zero (Sharpe 1.8 → 0.1), proving LLMs identify structural constraints independent of trading value—distinguishing algorithmic pattern recognition from the entrepreneurial discovery process that generates alpha.

Testing this framework on 242 trading days throughout 2024, we evaluate three manifestations of dealer hedging constraints: gamma positioning, stock pinning, and 0DTE hedging. Using fully unbiased prompts with obfuscated data, we achieve 71.5\% average detection rate, significantly exceeding the 60\% mechanical threshold. More importantly, 91.2\% of detected patterns materialize in forward returns, validating that the LLM identifies genuine market mechanics rather than spurious correlations.

These findings have important implications for both AI validation and financial market analysis. For the AI community, obfuscation testing offers a rigorous framework for validating reasoning capabilities in domains where training data contamination is a concern. For finance practitioners, our results suggest LLMs can serve as sophisticated pattern recognition tools for market microstructure analysis, provided their limitations are understood.

The remainder of this paper is organized as follows. Section II reviews related work. Section III presents our methodology. Section IV describes the experimental setup. Section V reports results. Section VI discusses implications and limitations. Section VII concludes.